\documentclass[12pt]{article}
\usepackage[utf8]{inputenc}
\usepackage[T1]{fontenc}
\usepackage{hyperref}
\title{Maore Comorian Language Resources for K-12 Curriculum Development}
\author{Compiled by Claude Sonnet 5.5 (claude-sonnet-5-5) for Andr\'as Kornai}
\date{}
\begin{document}
\maketitle

\section{Introduction}
Maore Comorian (Shimaore, Shimaor\'e) is spoken on Mayotte. Alphabets for Shimaore and Kibushi were adopted in 2020. A framework agreement signed on 8 June 2021 provides for training in, teaching of and promotion of the Mahorais languages. Teaching materials found online are scarce.

\section{Primary Educational Resources}
\begin{itemize}
\item Mayotte Hebdo, ``Vers un enseignement du shimaore et du kibushi \`a l'\'ecole ?'': \url{https://www.mayottehebdo.com/?p=3609}
\item Mayotte Hebdo, ``Enseignement des langues mahoraises \`a l'\'ecole, une nouvelle \'etape franchie'': \url{https://www.mayottehebdo.com/?p=26959}. It reports a convention allowing formal teaching of both languages in first-degree schools. No actual textbook was identified.
\item Written question and ministerial answer in the French National Assembly on regional languages in Mayotte schools: \url{https://www.assemblee-nationale.fr/dyn/15/questions/QANR5L15QE27356.pdf}. The answer notes that many teachers on the island speak neither language, so training is needed.
\end{itemize}

\section{Linguistic and Reference Materials}
\begin{itemize}
\item Lexilogos page on Comorian dictionaries, which links to an online Mahorais--French dictionary and a simplified grammar: \url{https://www.lexilogos.com/comorien_dictionnaire.htm}
\item Wikipedia entry on Shimaore, which cites Blanchy (1987), Cornice (1999) and Maandhui (1996): \url{https://wikipedia.com/wiki/Shimaore}
\end{itemize}

\section{Religious and Community Resources}
The association Shim\'e offers Shimaore and Kibushi classes in several communes. It is described in the Mayotte Hebdo articles above. I found no website for it. Mayotte Hebdo also reports on a mother-tongue celebration at Ouangani: \url{https://www.mayottehebdo.com/actualite/culture/les-langues-maternelles-celebrees-a-ouangani}

\section{Digital and Audio Resources}
No audio or video resources were found in these searches.

\section{Curriculum Development Context}
Formal teaching is only now beginning in primary schools, with teacher training as the main gap. Textbooks and graded readers could not be identified. A curriculum project would need to contact the Mayotte education authority and the association Shim\'e directly.

\end{document}